\section{Literature Review}

\subsection{SME Fragility and Resilience in Developing Economies}

The literature on SMEs in developing economies emphasizes both their economic importance and their vulnerability to institutional and market constraints. Earlier work shows that small firms account for a large share of private-sector activity but often face limited access to finance, weak infrastructure, and policy uncertainty \citep{beck2006small, ayyagari2007small}. Recent crisis-era evidence reinforces this point: firm resilience is shaped by liquidity, finance, technological capability, and the ability to reconfigure operations under disruption \citep{papadopoulos2020digital, guo2020digitalization, mishrif2023technology, amadasun2022finance}.

Research on firm resilience has increasingly moved beyond a simple failed/not-failed distinction. Organizational resilience studies view firms as having different capacities to anticipate, absorb, and adapt to shocks \citep{lengnickhall2005adaptive, duchek2020organizational}. In this framing, fragility can be interpreted as limited adaptive capacity under stress rather than as realized exit. Recent work on digital transformation and resilience similarly treats resilience as a dynamic capability built through learning, innovation, leadership, and technology-enabled adaptation \citep{awad2024digital, sagala2024toward, sagala2024antifragility, martinrojas2026resilience}.

\subsection{Digitalization, Innovation, and Modernization Risk}

Digitalization adds a further dimension. Firms that struggle to use digital tools may be less able to access customers, suppliers, payment systems, or information. Studies of SMEs during and after COVID-19 find that digital technologies supported crisis response, business continuity, and customer engagement, while also requiring complementary skills, finance, and organizational learning \citep{papadopoulos2020digital, guo2020digitalization, mishrif2023technology, awad2024digital, robertson2022digitalmaturity, lestari2024business, sinha2025digitalisation}. In emerging markets, this difficulty interacts with institutional voids such as uneven infrastructure, limited digital skills, and fragmented support systems \citep{khanna1997why}. The implication is not that digital presence or innovation is always protective. Modernization may also be reactive, costly, or incomplete when firms adopt new practices under pressure without the complementary resources needed to convert those practices into resilience.

\subsection{Predictive Modeling Versus Causal Inference in Firm-Risk Research}

Methodologically, many SME risk studies use descriptive statistics or conventional regression. Penalized regression offers a useful middle ground between interpretability and prediction. LASSO shrinks weak predictors toward zero and is especially useful when researchers want sparse models from many plausible covariates \citep{tibshirani1996regression, friedman2010regularization, hastie2015sparsity}. Stata and Python implementations now make regularized regression relatively accessible for applied social science, but model selection still depends on transparent cross-validation, preprocessing, and outcome construction \citep{ahrens2020lassopack, arlot2010crossvalidation}. In social-science applications, predictive performance should be interpreted carefully and separated from causal claims \citep{shmueli2010explain, mullainathan2017machine, athey2019machine}.

\subsection{Constructed Indices, Leakage, and Auditability}

The gap addressed by this paper is methodological as well as empirical. Development researchers often use constructed classifications because direct measures of closure, survival duration, or financial distress are unavailable. Such indices can be useful, but they require documentation of variable labels, coding direction, missing-value handling, and the relationship between index components and model predictors. When a model includes variables that are also used to construct the target outcome, high predictive accuracy may reflect index mechanics rather than independent empirical discovery \citep{kleinberg2018human}. This paper therefore contributes beyond applying LASSO to Rwanda Enterprise Survey data: it shows how to audit a constructed SME fragility index before treating it as evidence about firm distress or using it for policy screening.
